\section{Method}
\label{sec:method}

Let $p$ denote a candidate patch and $T = \{t_1, \ldots, t_n\}$ its test suite, where each $t_i$ has label $y_i \in \{\text{\textsc{Pass}}, \text{\textsc{Fail}}\}$ from execution. Tests decompose into \emph{fail-to-pass} ($T_{\text{f2p}}$) and \emph{pass-to-pass} ($T_{\text{p2p}}$) subsets; a patch is correct iff $\forall\, t_i{:}\ y_i = \text{\textsc{Pass}}$. We compare three formulations for predicting correctness without execution.

\subsection{Verification Formulations}
\label{sec:formulations}

Let $c_p$ denote the code context (unified diff) of patch $p$.

\paragraph{\cls{} (per-test classification).} $f(c_p, t_i) \rightarrow \hat{y}_i \in \{\text{\textsc{Pass}}, \text{\textsc{Fail}}\}$ predicts a single binary token per test (cross-entropy loss). Patch verdict: $\hat{y}_p = \bigwedge_{i} [\hat{y}_i = \text{\textsc{Pass}}]$. Output cost: one token per test.

\paragraph{\cwm{} (output generation).} $g(c_p, t_i) \rightarrow \hat{o}_i \in \Sigma^*$ generates the execution output string (language modeling loss). Verdict by exact match: $\hat{y}_i = \text{\textsc{Pass}}$ iff $\hat{o}_i = o_i^*$.\footnote{Exact string matching is conservative, consistent with \citet{meta2025cwm}. Semantic equivalences are not counted.} When the generated output is unparseable (i.e., does not conform to the expected structured format), the prediction defaults to the majority class verdict. Output cost: $\sum_i |\hat{o}_i|$ tokens.

\paragraph{\traj{} (trajectory-level).} $h(c_p, T) \rightarrow \hat{y}_p \in \{\text{\textsc{Resolved}}, \text{\textsc{Unresolved}}\}$ predicts a single patch verdict from all tests, providing no per-test signal. \traj{} serves as a granularity control: \cls{}--\traj{} differences isolate the value of per-test decomposition independent of the classification formulation itself.

\subsection{Formulation-Specific Failure Modes}
\label{sec:failure-modes}

Each formulation has characteristic failure modes that differ in their sensitivity to code complexity:

\paragraph{Generation failure modes (\cwm{}).} The open-ended output space introduces three failure modes: (1)~\emph{token-budget exhaustion}, where $|o_i^*| > B$ guarantees exact-match failure due to truncation; (2)~\emph{string-value prediction errors}, where semantically equivalent outputs differing in surface form (floating-point precision, dictionary ordering) produce false negatives; and (3)~\emph{format compliance collapse}, where the model fails to produce structured predictions entirely at high input complexity, even when the generation budget is sufficient. Under optimistic independence, a token sequence of length $L$ with per-token accuracy $q$ has correct-sequence probability $q^L$, making long outputs exponentially harder to predict exactly. We adopt exact string matching following the original \cwm{} evaluation protocol~\citep{meta2025cwm}. Unparseable predictions default to the majority-class label (\textsc{Pass}).

\paragraph{Classification failure modes (\cls{}).} The single-token output space eliminates format and truncation failures but constrains the model to encode its reasoning entirely within internal representations, without the intermediate reasoning trace that generation provides. \cls{} errors concentrate in false negatives (predicting \textsc{Pass} for failing tests), consistent with a bias toward the majority class.

\paragraph{Trajectory-level failure modes (\traj{}).} By collapsing all per-test information into a single binary verdict, \traj{} discards diagnostic granularity. When a patch fails multiple tests, the trajectory-level formulation cannot identify \emph{which} tests fail. This limits downstream utility (e.g., targeted re-repair) but avoids per-test aggregation errors.

\paragraph{Complexity dependence.} These failure modes suggest a complexity-dependent interaction: when execution outputs are short and structured (function-level code), generation modes (1)--(2) are rarely triggered, and all formulations should perform comparably. As output length and structural complexity increase (repository-level code), generation-specific failures may compound, creating divergence. Whether this divergence reflects a fundamental property of the generation formulation or is contingent on model/tokenizer properties is an empirical question we investigate in \S\ref{sec:experiments}.

\subsection{Oracle Accuracy-Utility Framework}
\label{sec:oracle}

Per-test granularity is valuable only when per-test predictions are sufficiently accurate and multiple tests provide partially independent signals. Consider a buggy patch with $k = |T_{\text{f2p}}|$ fail-to-pass tests. A per-test verifier with recall $r$ detects the bug if it identifies at least one failing test:
%
\begin{equation}
  P_{\text{detect}}(k) = 1 - (1 - r)^{k_{\text{eff}}}
  \label{eq:pdetect}
\end{equation}
%
where $k_{\text{eff}}$ is the effective number of independent tests. Let $a_h$ denote the trajectory-level detection accuracy; it is independent of $k$.

\paragraph{Block-correlation model.} Per-test predictions within the same patch are correlated. We model this via the intraclass correlation coefficient $\rho$ \citep[ICC;][]{shrout1979icc}:
%
\begin{equation}
  k_{\text{eff}} = \frac{k}{1 + (k - 1)\rho}
  \label{eq:keff}
\end{equation}
%
When $\rho = 0$, $k_{\text{eff}} = k$ (full independence); when $\rho = 1$, $k_{\text{eff}} = 1$ (per-test collapses to trajectory-level). The per-test advantage relative to a trajectory-level baseline is $\Delta(k) = P_{\text{detect}}(k) - a_h$, and the expected population advantage is:
%
\begin{equation}
  \mathbb{E}[\Delta] = \sum_{k=1}^{K} P(|T_{\text{f2p}}| = k) \cdot \Delta(k)
  \label{eq:expected-advantage}
\end{equation}

\paragraph{Testable predictions.} This yields three predictions validated in \S\ref{sec:experiments}: (1)~at $k = 1$, $\Delta \approx 0$; (2)~$\Delta(k)$ increases with $k$, dampened by $\rho > 0$; (3)~$\mathbb{E}[\Delta]$ concentrates in multi-F2P-test instances.
